\subsection{Operational Semantics of \GRSync}
\label{sec:semantics:grsync:opsem}

This section describes the operational semantics of the elements of \GRSync.
%
A state machine representation is common practice in \Lustre-like languages
\cite{DBLP:journals/tecs/ColacoMPP23}.
%
Similar to state machines, components compute output values based on current and past input values,
with states representing values needed for future computations: this includes values of the
identifiers $x$ referenced with \last{x} and states of invoked components.
%
The state has a bounded finite size, statically determined by the number of identifiers $x$ using
\last{x} and the number of called components.

\begin{table}[!ht]
    \centering
    \begin{tabularx}{\textwidth}{|l|X|}\hline
        $v, \, \some{v}, \, \none \in \Val$                             & A value of any type; \ttf{some} and \none represent optional values.                                                                                                                             \\\hline
        $v_\bot \in \ValBot ::= v \, \vert \, \bot$                     & A value that may be present ($v$) or absent ($\bot$).                                                                                                                                            \\\hline
        $[v^+]$                                                         & An implicitly finite list of values $v^+$.                                                                                                                                                       \\\hline
        $sty ::= ty \, \vert \, \mayB{ty}$                              & A type, where \kwf{?} indicates an event type.                                                                                                                                                   \\\hline
        $\tyCtx[x] = sty$                                               & A typing context \tyCtx mapping flow identifiers to their types.                                                                                                                                 \\\hline
        $\G[f]$                                                         & A program context \G mapping component and service names to their definitions.                                                                                                                   \\\hline
        $s$                                                             & The internal state of a component, an equation, or an expression.                                                                                                                                \\\hline
        \mkcell{$\rho[x] \, \vert \, \rho[\last{x}] = v_\bot$                                                                                                                                                                                                              \\
        $\rho ::= [(v_\bot/x)^+] \, \vert \, \rho_1 \ctxconcat \rho_2$} & A evaluation context $\rho$ mapping identifiers to their computed value.                                                                                                                         \\\hline
        \stateTy{\epsilon}                                              & The type of the compatible states of the \GRSync construct $\epsilon$.                                                                                                                           \\\hline
        $\initsync{\epsilon} \in \stateTy{\epsilon}$                    & The initial state of the \GRSync construct $\epsilon$.                                                                                                                                           \\\hline
        \compOpsem{s}{v^+_i}{f}{s'}{v^+_o}                              & Transition of the component $f$ with inputs $v_i^+$, updating its state and producing outputs $v_o^+$.                                                                                           \\\hline
        \eqOpsem{\rho}{s}{eq}{s'}{\rho'}                                & Transition of the equation $eq$ with a context $\rho$, updating its state and producing a new context $\rho'$.                                                                                   \\\hline
        \exprOpsem{\rho}{s}{e}{s'}{v^+_\bot}                            & Transition of the expression $e$ with a context $\rho$, updating its state and producing outputs $v^+_\bot$.                                                                                     \\\hline
        \patOpsem{\rho}{\guarded{pat^+ = v^+}{e_g}}{b_\bot}{\rho'}      & Transition from the context $\rho$ of the pattern $pat^+$ matching the values $v^+$, guarded by an optional expression $e_g$, and producing a bottom boolean $b_\bot$ and a new context $\rho'$. \\\hline
        \epatOpsem{\rho}{s}{\guarded{epat^+}{e_g}}{s'}{b_\bot}{\rho'}   & Transition from the context $\rho$ of the event pattern $epat^+$ guarded by an optional expression $e_g$, updating its state, and producing a bottom boolean $b_\bot$ and a new context $\rho'$. \\\hline
    \end{tabularx}
    \caption[Notations for the operational semantics of \GRSync.]%
    {Notations for the operational semantics of \GRSync.}
    \label{tab:semantics:grsync:opsem:notations}
\end{table}

We first introduce the notations used throughout this semantics, summarized in
\Cref{tab:semantics:grsync:opsem:notations}.
%
The set of all possible values is denoted by \Val, a value is denoted by $v$, and may be wrapped as
an optional using \some{v} or \none.
%
A bottom value $v_\bot$ can be either a defined value $v$ or the undefined marker $\bot$, used to
represent the presence or absence of a value. The set of all bottom values is then the addition of
\Val with the singleton $\{\bot\}$, denoted by \ValBot.
%
Across the entire semantics, we use the notation $X^+$ to denote a finite sequence of $X$, this
sequence can be empty.
%
Lists of elements $v^+$ are denoted by $[v^+]$.
%
The sampling types $sty$ of flows (\ie the type of the value visible in a flow at any given time)
are recorded in a typing context \tyCtx. A signal type is denoted by $ty$, while an event type is
denoted by \mayB{ty}, indicating that an event may have no value. The sampling type of a flow
identified by $x$ is accessed as $\tyCtx[x] = sty$. These types serve as syntactic labels, but their
correct usage is governed by a type system, detailed in \Cref{sec:grust:typing} of \Cref{cha:grust}.
In the semantics, they are used to distinguish the behavior of signals from that of events.
%
A program is represented as a mapping \G from component and service names to their definitions.
The definition of a component $f$ is accessed by $\G[F]$, and the one of a service $F$ by $\G[F]$.
%
States for components, equations, expressions are written $s$.
%
The evaluation context $\rho$ stores the computation state of the identifiers. For an identifier
$x$, $\rho[x] = \bot$ indicates that $x$ is not determined yet, and any operation involving this
bottom value also yields $\bot$. Otherwise, $\rho[x] = v$ gives the computed value for $x$.
%
The concatenation of two evaluation contexts $\rho_1$ and $\rho_2$, is defined as their disjoint
union, $\rho_1 \ctxconcat \rho_2$. An evaluation context mapping bottom values $v_\bot^+$ to
distinct identifiers $x^+$ is denoted by $[(v_\bot/x)^+]$.
%
The execution of \GRSync constructs is modeled as a state machine.
For all \GRSync construct $\epsilon$, the set of all states corresponding to $\epsilon$ is denoted
by \stateTy{\epsilon}, and $\initsync{\epsilon} \in \stateTy{\epsilon}$ denotes its initial state.
The following paragraphs explain in details the computation steps of the different constructs.

\paragraph{Components}

The transition function of a component $f$, defined in a program \G with a typing context \tyCtx,
is denoted by the judgement
\begin{equation*}
    \compOpsem{s}{v_i^+}{f}{s'}{v_o^+}.
\end{equation*}
%
From a state $s$ and with input values $v^+_i$, it produces the next state $s'$ with output values
$v^+_o$.

The initial state of a component $f$ defined in a program \G is denoted by $\initsync{f} \in \stateTy{f}$ and
defined in \Cref{eq:semantics:grsync:opsem:init:comp}. It corresponds to the initialization of the
state of its internal equations and memories.
\begin{align}
    \stateTy{f}  & = \Pi_{x_m} \Val \times \stateTy{eq^+} \label{eq:semantics:grsync:opsem:stateTy:comp} \\
    \initsync{f} & = (c^+, \initsync{eq^+}) \label{eq:semantics:grsync:opsem:init:comp}                  \\
                 & \text{with} \, G[f] = \component{f}{x_i^+}{x_o^+}{eq^+}{x_m^+ = c^+} \nonumber
\end{align}

As defined by rule \nameref{rule:semantics:grsync:opsem:step:comp}, the transition function of a
component performs the transition of its internal equations ($eq^+$) in a context initialized with
the input values ($(v_i/x_i)^+$) and the memory values ($(v_m/\last{x_m})^+$) retrieved from its
internal state ($(v_m^+, s)$).

\begin{center}
    \begin{minipage}{0.9\textwidth}
        \centering
        \begin{mathpar}
            %%%%% Component
            \inferOPSEM
            {
            G[f] = \component{f}{x_i^+}{x_o^+}{eq^+}{x_m^+ = c^+} \\\\
            \eqOpsem{[(v_i/x_i)^+, (v_m/\last{x_m})^+]}{s}{eq^+}{s'}{\rho} \\
            \left(v_m' = \rho[x_m]\right)^+ \\
            \left(v_o = \rho[x_o]\right)^+
            }
            { \compOpsem{(v_m^+, s)}{v_i^+}{f}{(v_m^{\prime+}, s')}{v_o^+} }
            {Comp}
            \label{rule:semantics:grsync:opsem:step:comp}
        \end{mathpar}
    \end{minipage}
\end{center}

For instance, in the \aeb example~\Cref{fig:semantics:example:aeb}, the initial state and transition function
of the component \idf{braking\_strat} are $\initsync{\idf{braking\_strat}} = \No$ and
\opsem{S, I}{\idf{braking\_strat}}{S', O}.
Starting from the state $S = \No$ and inputs $I = \{\idf{s} = 32., \idf{p} = \some{20.},
    \idf{t\_p}=\none\}$
it produces a new state $S' = \Soft$ and outputs $O = \{\idf{b}=\Soft\}$.

\paragraph{Equations}

While equations don't have direct input values, they rely on the evaluation context $\rho$ to access
the values of the already evaluated identifiers.
%
Their initial state is denoted by  $\initsync{eq} \in \stateTy{eq}$, and their transition in an
evaluation context $\rho$ by
\begin{equation*}
    \eqOpsem{\rho}{s}{eq}{s'}{\rho'}.
\end{equation*}
Starting from state a $s$, it produces a new state $s'$ and a new context $\rho'$, which includes
the identifiers defined by the equations.

\subparagraph{Blocks of equations}

The equations in a block ($eq^+$) are simultaneous, all must be computed in the same transition.
But an equation may depend on identifiers defined by other equations, so a causality check
(formalized in \Cref{def:semantics:grsync:causality:causal:prop}) verifies there exists an order in
which equations can be computed without any deadlock caused by circular dependencies
(this scheduling property is stated in \Cref{prop:semantics:grsync:causality:scheduling}).

The initial state of a block of equations $eq^+$ in a program \G is denoted by \initsync{eq^+}
and defined in \Cref{eq:semantics:grsync:opsem:init:eqs}. It corresponds to the initialization of
all individual equations.
\begin{align}
    \stateTy{eq^+}  & = \Pi_{eq} \stateTy{eq} \label{eq:semantics:grsync:opsem:stateTy:eqs} \\
    \initsync{eq^+} & = (\initsync{eq})^+ \label{eq:semantics:grsync:opsem:init:eqs}
\end{align}

\begin{center}
    \begin{minipage}{0.9\textwidth}
        \centering
        \begin{mathpar}
            %%%%% Block of equations
            \inferOPSEMEq
            { \fixsem{n+1}{\rho}{(s_1, \dots, s_n)}{eq_1, \dots, eq_n}{(s'_1, \dots, s'_n)}{\rho'} }
            { \eqOpsem{\rho}{(s_1, \dots, s_n)}{eq_1, \dots, eq_n}{(s'_1, \dots, s'_n)}{\rho'} }
            {Blck}
            \label{rule:semantics:grsync:opsem:step:eqs}

            %%%%% Fix-point
            \inferOPSEMEq
            { }
            { \fixsem{0}{\rho}{s^+}{eq^+}{s^+}{\botrho{\defs{eq^+}}} }
            {Fix{\scriptsize 0}}
            \label{rule:semantics:grsync:opsem:step:fix0}

            \inferOPSEMEq
            {
            \fixsem{k}{\rho}{s_i^+}{eq_i^+}{s_i^{\prime\prime +}}{\rho'} \\\\
            \forall i, \, \eqOpsem{\rho \ctxconcat \rho'}{s_i}{eq_i}{s'_i}{\rho'_i}
            }
            { \fixsem{k+1}{\rho}{s_i^+}{eq_i^+}{s_i^{\prime +}}{\bigCtxconcat_i \rho'_i} }
            {Fix{\scriptsize k+1}}
            \label{rule:semantics:grsync:opsem:step:fixkp1}
        \end{mathpar}
    \end{minipage}
\end{center}

The transition of $eq^+$, described by rule \nameref{rule:semantics:grsync:opsem:step:eqs},
uses the fix-point transition, defined by the rules \nameref{rule:semantics:grsync:opsem:step:fix0}
and \nameref{rule:semantics:grsync:opsem:step:fixkp1}. The latter loops on the transition of all
individual equations $eq$, computing iteratively the undetermined identifiers using the values of
the already determined ones.
%
During a transition, if the value of an identifier is not determined yet, then attempting to
retrieve it from the context $\rho$ returns $\bot$, and \emph{any expression involving this bottom
    value also yields $\bot$} (as shown in the paragraph on expressions).
%
The fix-point transition starts with the empty context \botrho{\defs{eq^+}} that associates
$\bot$ to all the identifiers defined by the equations
(\defs{eq^+} is defined in \Cref{sec:semantics:grsync:causality:defs}).

The maximum number of iterations is equal to the number of simultaneous equations. Intuitively, each
iteration resolves at least one previously unresolved equation; if no new equations are resolved
during an iteration, the context $\rho$ remains unchanged, indicating that a fix-point has been
reached. This point will be formally proved in \Cref{thm:semantics:grsync:theorems:productivity}.

\begin{remark}
    \label{rmk:semantics:grsync:opsem:step:eqs}
    We adopt a least fix-point semantics instead of relying on a static scheduling of the
    equations based on their dependency graph, as performed by the compiler.
    This choice reflects our goal to describe the operational semantics of the program itself,
    rather than an intermediate compilation stage. Moreover, non-static scheduling enables
    concurrent evaluation of independent equations, opening the door to parallel execution
    strategies.
\end{remark}

\subparagraph{Declarations}

The transition of a declaration $x^+ = e$ is defined by the rule
\nameref{rule:semantics:grsync:opsem:step:eq:decl} that handles the transition of the expression $e$
and defines the context $[(v_\bot / x)^+]$, where the identifiers $x^+$ are mapped to the
bottom values~$v_\bot^+$ obtained from the transition of the expression $e$.

The initial state of declarations corresponds to the initial state of their expressions, as stated
by \Cref{eq:semantics:grsync:opsem:init:eq:decl}.

\begin{align}
    \stateTy{x^+ = e} & = \stateTy{e} & \initsync{x^+ = e} & = \initsync{e} \label{eq:semantics:grsync:opsem:init:eq:decl}
\end{align}

\begin{center}
    \begin{minipage}{0.45\textwidth}
        \centering
        \begin{mathpar}
            %%%%% Declarations
            \inferOPSEMEq
            { \exprOpsem{\rho}{s}{e}{s'}{v_\bot^+} }
            { \eqOpsem{\rho}{s}{x^+ = e}{s'}{[(v_\bot / x)^+]} }
            {Decl}
            \label{rule:semantics:grsync:opsem:step:eq:decl}
        \end{mathpar}
    \end{minipage}
\end{center}

\subparagraph{Match}

The transition of \kwf{match} equations is described by the rules
\nameref{rule:semantics:grsync:opsem:step:eq:match:value} and
\nameref{rule:semantics:grsync:opsem:step:eq:match:bot}.

It performs the transition of the matched expression and, if it returns values, following
\nameref{rule:semantics:grsync:opsem:step:eq:match:value}, uses those values to test the guarded
patterns of all branches. Each of those patterns returns a boolean telling whether the guarded
pattern matches or not the value, and a context that contains identifiers locally introduced by the
pattern. Then it evaluates the equations of the first pattern $pat_k^+$ to match the values of $e$
with a guard $e_k$ that holds. Exhaustive patterns ensure that this transition is always possible.

If the matched expression returns one bottom value, following the rule
\nameref{rule:semantics:grsync:opsem:step:eq:match:bot}, then the pattern matching cannot continue
further, returning the context \botrho{\{x^+\}} that maps every defined identifiers to $\bot$.

The initial state of \kwf{match} equations corresponds to the initial state of the matched
expression $e$ and the initial states of branches equations, as stated by
\Cref{eq:semantics:grsync:opsem:init:eq:match}. The guards are memory free, no states are required
for them.

\begin{align}
    \stateTy{\mathfoot{\matchTab{e}{x^+}{(\branch{pat^+}{e_g}{eq^+})^+}}}  & = \stateTy{e} \times \Pi_{eq^+} \stateTy{eq^+}
    \label{eq:semantics:grsync:opsem:stateTy:eq:match}                                                                      \\
    \initsync{\mathfoot{\matchTab{e}{x^+}{(\branch{pat^+}{e_g}{eq^+})^+}}} & = (\initsync{e}, (\initsync{eq^+})^+)
    \label{eq:semantics:grsync:opsem:init:eq:match}
\end{align}

\begin{center}
    \begin{minipage}{0.99\textwidth}
        \centering
        \begin{mathpar}
            %%%%% Match equation
            \inferOPSEMEq
            {
            \exprOpsem{\rho}{s_e}{e}{s'_e}{v_e^+} \\
            \forall i, \, \patOpsem{\rho}{\guarded{pat_i^+ = v_e^+}{e_i}}{b_i}{\rho_i} \\\\
            \exists k, \, b_k = \true \\
            \forall i < k, \, b_i = \false \\\\
            \eqOpsem{\rho \ctxconcat \rho_k}{s_k}{eq_k^+}{s'_k}{\rho'} \\
            \forall i \neq k, \, s'_i = s_i
            }
            { \eqOpsem{\rho}{(s_e, s_i^+)}{\match{e}{x^+}{(\branch{pat_i^+}{e_i}{eq_i^+})^+}
            }{(s'_e, s_i^{\prime+})}{\rho'} }
            {Mtch{\scriptsize val}}
            \label{rule:semantics:grsync:opsem:step:eq:match:value}

            \inferOPSEMEq
            {
                \exprOpsem{\rho}{s_e}{e}{s'_e}{v_{\bot e}^+} \\
                \exists v_\bot \in \{v_{\bot e}^+\}, \: v_\bot = \bot
            }
            { \eqOpsem{\rho}{(s_e, s_i^+)}{\match{e}{x^+}{(\branch{pat_i^+}{e_i}{eq_i^+})^+}
                }{(s_e, s_i^+)}{\botrho{\{x^+\}}} }
            {Mtch{\scriptsize bot}}
            \label{rule:semantics:grsync:opsem:step:eq:match:bot}
        \end{mathpar}
    \end{minipage}
\end{center}

\subparagraph{When}

The semantics of \kwf{when} equations is similar to \kwf{match} equations. It performs the
transition of the guarded event patterns of each branches, which returns a boolean telling whether
the guarded event pattern is triggered or not, and a context that contains identifiers locally
introduced by the event pattern.

If all pattern returns a boolean (not bottom) and if there is a branch whose event pattern returns
\true, the transition follows the rule \nameref{rule:semantics:grsync:opsem:step:eq:when:branch}. It
performs the updates defined by the equations $eq_k^+$ of the first event pattern $epat_k^+$ that is
triggered with a guard $e_k$ that holds.
%
Identifiers not in $eq^+_k$ are added to the context evaluation context of $eq_k^+$ with default
values in $\rho_k'$: previous value for signals (identifiers that are initialized in
\init{x_m^+ = c^+}), \none for events.

If all pattern returns a boolean but no event pattern is triggered, the rule
\nameref{rule:semantics:grsync:opsem:step:eq:when:default} applies and all identifiers are set to
default.

If there is a pattern that returns $\bot$, following
\nameref{rule:semantics:grsync:opsem:step:eq:when:bot} the \kwf{when} equation cannot continue
further, returning the context \botrho{\{x^+\}} that maps every defined identifiers to $\bot$.

Initial states for \kwf{when} equations contain the initial values of the signals it defines, the
initial states of event patterns and branches equations, as stated by
\Cref{eq:semantics:grsync:opsem:init:eq:when}. The guards are also memory free.

\begin{align}
    \stateTy{\mathfoot{\whenTab{x^+}{x_m^+ = c}{(\branch{epat^+}{e_g}{eq^+})^+}}}  & =
    \begin{array}{l}
        \Pi_{x_m} \Val \times \Pi_{epat^+} \stateTy{epat^+} \\
        \times \Pi_{eq^+} \stateTy{eq^+}
    \end{array}
    \label{eq:semantics:grsync:opsem:stateTy:eq:when}                                  \\
    \initsync{\mathfoot{\whenTab{x^+}{x_m^+ = c}{(\branch{epat^+}{e_g}{eq^+})^+}}} & =
    (c^+, (\initsync{epat^+})^+ (\initsync{eq^+})^+)
    \label{eq:semantics:grsync:opsem:init:eq:when}
\end{align}

\begin{center}
    \begin{minipage}{0.99\textwidth}
        \centering
        \begin{mathpar}
            %%%%% When equation
            \inferOPSEMEq
            {
            \forall i, \, \epatOpsem{\rho}{s^p_i}{\guarded{epat_i^+}{e_i}}{s^{p\prime}_i}{b_i}{\rho_i} \\
            \exists k, \, b_k = \true \\
            \forall i < k, \, b_i = \false \\\\
            %
            E_k = (\{x^+\} \backslash \{x_m^+\}) \backslash \defs{eq^+_k} \\
            S_k = \{x_m^+\} \backslash \defs{eq^+_k} \\
            %
            \rho_m = [(v_m / \last{x_m})^+] \\
            \rho_k' = [\none / x_e \in E_k] \ctxconcat [\rho_m[\last{x_s}] / x_s \in S_k] \\
            %
            \eqOpsem{\rho \ctxconcat \rho_k \ctxconcat \rho_m \ctxconcat \rho_k'}{s_k}{eq_k^+}{s'_k}{\rho'} \\
            %
            \left(v_m' = (\rho'_k \ctxconcat \rho')[x_m]\right)^+ \\
            \forall i \neq k, \, s'_i = s_i
            }
            { \eqOpsem{\rho}{(v_m^+, s_i^{p+}, s_i^+)}{
            \whenTab{x^+}{x_m^+ = c^+}{(\branch{epat_i^+}{e_i}{eq_i^+})^+}
            }{(v_m^{\prime+}, s_i^{p\prime+}, s_i^{\prime+})}{\rho_m \ctxconcat \rho'_k \ctxconcat \rho'} }
            {When{\scriptsize trig}}
            \label{rule:semantics:grsync:opsem:step:eq:when:branch}

            \inferOPSEMEq
            {
            \forall i, \, \epatOpsem{\rho}{s^p_i}{\guarded{epat_i^+}{e_i}}{s^{p\prime}_i}{\false}{\rho_i} \\\\
            \rho_m = [(v_m / \last{x_m})^+] \\
            \rho' = \rho_m \ctxconcat [\none / x_e \in \{x^+\} \backslash \{x_m^+\}] \ctxconcat [(v_m / x_m)^+]
            }
            { \eqOpsem{\rho}{(v_m^+, s_i^{p+}, s_i^+)}{
            \whenTab{x^+}{x_m^+ = c^+}{(\branch{epat_i^+}{e_i}{eq_i^+})^+}
            }{(v_m^+, s_i^{p\prime+}, s_i^+)}{\rho'} }
            {When{\scriptsize none}}
            \label{rule:semantics:grsync:opsem:step:eq:when:default}

            \inferOPSEMEq
            { \exists i, \, \epatOpsem{\rho}{s^p_i}{\guarded{epat_i^+}{e_i}}{s^{p\prime}_i}{\bot}{\rho_i} }
            { \eqOpsem{\rho}{(v_m^+, s_i^{p+}, s_i^+)}{
            \whenTab{x^+}{x_m^+ = c^+}{(\branch{epat_i^+}{e_i}{eq_i^+})^+}
            }{(v_m^+, s_i^{p+}, s_i^+)}{\botrho{\{x^+\}}} }
            {When{\scriptsize bot}}
            \label{rule:semantics:grsync:opsem:step:eq:when:bot}
        \end{mathpar}
    \end{minipage}
\end{center}

For instance, in the \aeb example~\Cref{fig:semantics:example:aeb}, the initial state of the \kwf{when}
equation at \Crefrange{ex:line:whenStart}{ex:line:whenEnd} is explicitly defined in
\Cref{eq:semantics:grsync:opsem:init:eq:when:ex}. As a matter of fact, the state of the component
\idf{compute\_braking} is empty, so we can simplify the state of the \kwf{when} equation to the
memory of \idf{b} that is initialized to \No.
\begin{equation}
    \initsync{\mathtiny{\whenTab{\idf{b}}{\idf{b} = \No}{
                \begin{array}{l}
                    \detect{\idf{x}}{\idf{p}} \Rightarrow { \idf{b} = \idf{compute\_braking}(\idf{x}, \idf{s}); } \\
                    \detect{\idf{x}}{\idf{t\_p}} \Rightarrow { \idf{b} = \No;  }                                  \\
                \end{array}
            }}}
    \begin{array}{l}
        = (\No, (), (\initsync{\idf{compute\_braking}}, ())) \\
        = (\No, (), ((), ()))                                \\
        \sim \No
    \end{array}
    \label{eq:semantics:grsync:opsem:init:eq:when:ex}
\end{equation}

Based on the chronogram from \Cref{fig:semantics:example:aeb:chronogram}, we can model the execution
of the \kwf{when} equation into the sequence of transitions from
\Crefrange{eq:semantics:grsync:opsem:step:eq:when:ex:1}{eq:semantics:grsync:opsem:step:eq:when:ex:5}.
In \Cref{%
    eq:semantics:grsync:opsem:step:eq:when:ex:1,%
    eq:semantics:grsync:opsem:step:eq:when:ex:3,%
    eq:semantics:grsync:opsem:step:eq:when:ex:4%
}, all the events are absent and set to \none in the input evaluation context, so the rule
\nameref{rule:semantics:grsync:opsem:step:eq:when:default} is applied and the value of \idf{b}
remains the same. The event \idf{p} is present in
\Cref{eq:semantics:grsync:opsem:step:eq:when:ex:2}, so the rule
\nameref{rule:semantics:grsync:opsem:step:eq:when:branch} is applied with the first branch being
triggered, such that the value of \idf{b} and the state of the \kwf{when} equation are updated to
\Soft. The timeout event \idf{t\_p} is present in
\Cref{eq:semantics:grsync:opsem:step:eq:when:ex:5}, so the rule
\nameref{rule:semantics:grsync:opsem:step:eq:when:branch} is applied with the second branch being
triggered, such that the value of \idf{b} and the state of the \kwf{when} equation are reset to \No.
\begin{align}
    \eqOpsem
    {[32. / \idf{s}, \none / \idf{p}, \none / \idf{t\_p}]}
    {\No   & }
    {\kwf{when} \{ \dots \}}
    {\No}
    {[\No / \idf{b}, \No / \last{\idf{b}}]}
    \label{eq:semantics:grsync:opsem:step:eq:when:ex:1} \\
    %
    \eqOpsem
    {[32. / \idf{s}, \some{20.} / \idf{p}, \none / \idf{t\_p}]}
    {\No   & }
    {\kwf{when} \{ \dots \}}
    {\Soft}
    {[\Soft / \idf{b}, \No / \last{\idf{b}}]}
    \label{eq:semantics:grsync:opsem:step:eq:when:ex:2} \\
    %
    \eqOpsem
    {[30. / \idf{s}, \none / \idf{p}, \none / \idf{t\_p}]}
    {\Soft & }
    {\kwf{when} \{ \dots \}}
    {\Soft}
    {[\Soft / \idf{b}, \Soft / \last{\idf{b}}]}
    \label{eq:semantics:grsync:opsem:step:eq:when:ex:3} \\
    %
    \eqOpsem
    {[25. / \idf{s}, \none / \idf{p}, \none / \idf{t\_p}]}
    {\Soft & }
    {\kwf{when} \{ \dots \}}
    {\Soft}
    {[\Soft / \idf{b}, \Soft / \last{\idf{b}}]}
    \label{eq:semantics:grsync:opsem:step:eq:when:ex:4} \\
    %
    \eqOpsem
    {[25. / \idf{s}, \none / \idf{p}, \some{()} / \idf{t\_p}]}
    {\Soft & }
    {\kwf{when} \{ \dots \}}
    {\No}
    {[\No / \idf{b}, \Soft / \last{\idf{b}}]}
    \label{eq:semantics:grsync:opsem:step:eq:when:ex:5}
\end{align}

\paragraph{Expressions}

Similarly to equations, \GRSync expressions don't have input values but an input evaluation context
$\rho$ that contains the computed identifiers.
%
Their initial state is denoted by  $\initsync{e} \in \stateTy{e}$, and their transition with an evaluation context
$\rho$ by
\begin{equation*}
    \exprOpsem{\rho}{s}{e}{s'}{v_\bot^+}.
\end{equation*}
Starting from state a $s$, it produces a new state $s'$ and output bottom values $v_\bot^+$.

The rules for constant \nameref{rule:semantics:grsync:opsem:step:e:const} and identifier
\nameref{rule:semantics:grsync:opsem:step:e:ident} expressions use an empty state
(\Cref{eq:semantics:grsync:opsem:init:e:ident,eq:semantics:grsync:opsem:init:e:const}).
They produce respectively the constant value and the value stored in the context $\rho$.

\begin{multicols}{2}[]
    \begin{align}
        \stateTy{c} & = \{()\} & \initsync{c} & = () \label{eq:semantics:grsync:opsem:init:e:const}
    \end{align}

    \newcolumn

    \begin{minipage}{0.45\textwidth}
        \begin{mathpar}
            \inferOPSEME
            { }
            { \exprOpsem{\rho}{()}{c}{()}{c} }
            {Cst}
            \label{rule:semantics:grsync:opsem:step:e:const}
        \end{mathpar}
    \end{minipage}
\end{multicols}
\begin{multicols}{2}[]
    \begin{align}
        \stateTy{x} & = \{()\} & \initsync{x} & = () \label{eq:semantics:grsync:opsem:init:e:ident}
    \end{align}

    \newcolumn

    \begin{minipage}{0.45\textwidth}
        \begin{mathpar}
            \inferOPSEME
            { }
            { \exprOpsem{\rho}{()}{x}{()}{\rho[x]} }
            {ID}
            \label{rule:semantics:grsync:opsem:step:e:ident}
        \end{mathpar}
    \end{minipage}
\end{multicols}

The state of a memory expression (\last{x}) is also empty
(\Cref{eq:semantics:grsync:opsem:init:e:last}), because the last value of $x$ is stored in the
context $\rho$, managed by initializations (\init{x_m^+ = c^+}).
The transition of (\last{x}) only retrieves this memory, as depicted by rule
\nameref{rule:semantics:grsync:opsem:step:e:last}.

\begin{multicols}{2}[]
    \begin{align}
        \stateTy{\last{x}} & = \{()\} & \initsync{\last{x}} & = () \label{eq:semantics:grsync:opsem:init:e:last}
    \end{align}

    \newcolumn

    \begin{minipage}{0.45\textwidth}
        \begin{mathpar}
            \inferOPSEME
            { }
            { \exprOpsem{\rho}{()}{\last{x}}{()}{\rho[\last{x}]} }
            {Last}
            \label{rule:semantics:grsync:opsem:step:e:last}
        \end{mathpar}
    \end{minipage}
\end{multicols}

The initial state of an event emission (\emit{e}) only contains the initial state of the expression
emitted ($e$), as stated by \Cref{eq:semantics:grsync:opsem:init:e:emit}. Its transition, as defined
by rule \nameref{rule:semantics:grsync:opsem:step:e:emit:value}, wraps the value returned by the
transition of $e$ in a \ttf{some} option to indicate the event is present. If the transition of $e$
returns a bottom value, then \emit{e} follows rule
\nameref{rule:semantics:grsync:opsem:step:e:emit:bot} and returns $\bot$ as well, indicating that
the computation cannot yet be performed.

\begin{align}
    \stateTy{\emit{e}} & = \stateTy{e} & \initsync{\emit{e}} & = \initsync{e} \label{eq:semantics:grsync:opsem:init:e:emit}
\end{align}

\begin{center}
    \begin{minipage}{0.9\textwidth}
        \begin{mathpar}
            \inferOPSEME
            { \exprOpsem{\rho}{s}{e}{s'}{v} }
            { \exprOpsem{\rho}{s}{\emit{e}}{s'}{\some{v}} }
            {Emit{\scriptsize val}}
            \label{rule:semantics:grsync:opsem:step:e:emit:value}

            \inferOPSEME
            { \exprOpsem{\rho}{s}{e}{s'}{\bot} }
            { \exprOpsem{\rho}{s}{\emit{e}}{s}{\bot} }
            {Emit{\scriptsize bot}}
            \label{rule:semantics:grsync:opsem:step:e:emit:bot}
        \end{mathpar}
    \end{minipage}
\end{center}

The initial states of $n$-ary operations ($op(e^+)$), defined by
\Cref{eq:semantics:grsync:opsem:init:e:op}, contain the initial states of the input expressions
($e^+$). Their transition function, defined in rule
\nameref{rule:semantics:grsync:opsem:step:e:op:value}, evaluates the transitions of the input
expressions $e^+$ and performs the operation on their returned values. If the transition of one
input expression returns a bottom value, then $op(e^+)$ follows rule
\nameref{rule:semantics:grsync:opsem:step:e:op:bot} and returns $\bot$ as well, indicating that
the computation cannot yet be performed.

\begin{align}
    \stateTy{op(e^+)}  & = \Pi_e \stateTy{e}                                            &
    \initsync{op(e^+)} & = (\initsync{e})^+ \label{eq:semantics:grsync:opsem:init:e:op}
\end{align}

\begin{center}
    \begin{minipage}{0.95\textwidth}
        \begin{mathpar}
            \inferOPSEME
            { \left(\exprOpsem{\rho}{s}{e}{s'}{v}\right)^+ }
            { \exprOpsem{\rho}{s^+}{op(e^+)}{s^{\prime+}}{op(v^+)} }
            {Op{\scriptsize val}}
            \label{rule:semantics:grsync:opsem:step:e:op:value}

            \inferOPSEME
            {
                \left(\exprOpsem{\rho}{s}{e}{s'}{v_{\bot i}}\right)^+ \\\\
                \exists v_\bot \in \{v_{\bot i}^+\}, \; v_\bot = \bot
            }
            { \exprOpsem{\rho}{s^+}{op(e^+)}{s^+}{\bot} }
            {Op{\scriptsize bot}}
            \label{rule:semantics:grsync:opsem:step:e:op:bot}
        \end{mathpar}
    \end{minipage}
\end{center}

The initial states of component applications ($f(e^+)$), as stated by
\Cref{eq:semantics:grsync:opsem:init:e:app}, contain the initial states of the input expressions
($e^+$) plus the initial state of the component ($f$). Thus, only the states of called components
are stored in expressions memory. For component applications, rule
\nameref{rule:semantics:grsync:opsem:step:e:app:value} applies: it evaluates the transitions of the
input expressions $e^+$, and uses their values as inputs in the transition the component $f$. If the
transition of one input expression returns a bottom value, then $f(e^+)$ follows rule
\nameref{rule:semantics:grsync:opsem:step:e:app:bot} and returns $\bot$ as well, indicating that
the computation cannot yet be performed.

\begin{align}
    \stateTy{f(e^+)}  & = \stateTy{f} \times \Pi_e \stateTy{e}                                          &
    \initsync{f(e^+)} & = (\initsync{f}, (\initsync{e})^+) \label{eq:semantics:grsync:opsem:init:e:app}
\end{align}

\begin{center}
    \begin{minipage}{0.95\textwidth}
        \begin{mathpar}
            \inferOPSEME
            {
                \left(\exprOpsem{\rho}{s}{e}{s'}{v_i}\right)^+ \\
                \compOpsem{s_f}{v_i^+}{f}{s'_f}{v_o^+}
            }
            { \exprOpsem{\rho}{(s_f, s^+)}{f(e^+)}{(s'_f, s^{\prime+})}{v_o^+} }
            {App{\scriptsize val}}
            \label{rule:semantics:grsync:opsem:step:e:app:value}

            \inferOPSEME
            {
                \left(\exprOpsem{\rho}{s}{e}{s'}{v_{\bot i}}\right)^+ \\
                \exists v_\bot \in \{v_{\bot i}^+\}, \; v_\bot = \bot
            }
            { \exprOpsem{\rho}{(s_f, s^+)}{f(e^+)}{(s_f, s^+)}{\bot} }
            {App{\scriptsize bot}}
            \label{rule:semantics:grsync:opsem:step:e:app:bot}
        \end{mathpar}
    \end{minipage}
\end{center}

\paragraph{Patterns}

The execution of \GRSync patterns indicates if the pattern holds or not, and defines some local
identifiers.
%
Their states are always empty because their guards must memory free, as stated by
\Cref{eq:semantics:grsync:opsem:stateTy:pats,eq:semantics:grsync:opsem:init:pats}.
\begin{align}
    \stateTy{\guarded{pat^+}{e_g}}  & = \{()\} & \text{\iff } & \stateTy{e_g} = \{()\} \label{eq:semantics:grsync:opsem:stateTy:pats} \\
    \initsync{\guarded{pat^+}{e_g}} & = ()     & \text{\iff } & \initsync{e_g} = () \label{eq:semantics:grsync:opsem:init:pats}
\end{align}
The transition of a guarded pattern \guarded{pat^+}{e_g}, matching the values $v^+$ with an
evaluation context $\rho$, produces an output bottom boolean $b_\bot$ that tells if the pattern
matches the values and the guard holds ($\bot$ means that matching cannot yet be determined), and
produces a new evaluation context $\rho'$ that contains the identifiers locally defined by the
pattern. This transition is denoted by:
\begin{equation*}
    \patOpsem{\rho}{\guarded{pat^+ = v^+}{e_g}}{b_\bot}{\rho'}.
\end{equation*}

The transition for guarded multiple patterns matching some values (\guarded{pat^+ = v^+}{e_g}),
ruled by \nameref{rule:semantics:grsync:opsem:step:pats:bool}, executes the transitions of all
individual patterns matching the corresponding value ($pat = v$) and the transition of the guard
$e_g$. The pattern holds only if all individual patterns and the guard hold. The locally defined
identifiers are combined in a one evaluation context. If the guard returns $\bot$, then the rule
\nameref{rule:semantics:grsync:opsem:step:pats:bot} applies and the pattern returns $\bot$ as well.

\begin{center}
    \begin{minipage}{0.9\textwidth}
        \begin{mathpar}
            \inferOPSEMPat
            {
                \left(\patOpsem{\rho}{pat = v}{b'}{\rho'}\right)^+  \\\\
                \exprOpsem{\rho \ctxconcat \bigCtxconcat \rho^{\prime+}}{()}{e_g}{()}{b_g} \\
                b = b_g \kwf{\&\&} (\kwf{\&\&} b')^+
            }
            { \patOpsem{\rho}{\guarded{pat^+ = v^+}{e_g}}{b}{\bigCtxconcat \rho^{\prime +}} }
            {Tple{\scriptsize bool}}
            \label{rule:semantics:grsync:opsem:step:pats:bool}

            \inferOPSEMPat
            {
                \left(\patOpsem{\rho}{pat = v}{b'}{\rho'}\right)^+  \\\\
                \exprOpsem{\rho \ctxconcat \bigCtxconcat \rho^{\prime+}}{()}{e_g}{()}{\bot}
            }
            { \patOpsem{\rho}{\guarded{pat^+ = v^+}{e_g}}{\bot}{\bigCtxconcat \rho^{\prime +}} }
            {Tple{\scriptsize bot}}
            \label{rule:semantics:grsync:opsem:step:pats:bot}
        \end{mathpar}
    \end{minipage}
\end{center}

The transition of wildcards ($\_ = v$) and identifier patterns ($x = v$) always hold, as the rules
\nameref{rule:semantics:grsync:opsem:step:pat:wild} and
\nameref{rule:semantics:grsync:opsem:step:pat:ident} state. There is no locally defined identifier
for wildcards, but the identifier pattern $x = v$ defines the identifier $x$ to be associated with the
matching value $v$.

\begin{center}
    \begin{minipage}{0.9\textwidth}
        \begin{mathpar}
            \inferOPSEMPat
            { }
            { \patOpsem{\rho}{\_ = v}{\true}{\emptyrho} }
            {Wild}
            \label{rule:semantics:grsync:opsem:step:pat:wild}

            \inferOPSEMPat
            { }
            { \patOpsem{\rho}{x = v}{\true}{[v / x]} }
            {ID}
            \label{rule:semantics:grsync:opsem:step:pat:ident}
        \end{mathpar}
    \end{minipage}
\end{center}

The transition of a constant pattern matching $c = v$ is defined by the rule
\nameref{rule:semantics:grsync:opsem:step:pat:const}: it states that the pattern holds if and only
if the constant $c$ and the value $v$ are equal. No identifiers are locally defined by this pattern.

\begin{center}
    \begin{minipage}{0.9\textwidth}
        \begin{mathpar}
            \inferOPSEMPat
            { }
            { \patOpsem{\rho}{c = v}{c \kwf{==} v}{\emptyrho} }
            {Cst}
            \label{rule:semantics:grsync:opsem:step:pat:const}
        \end{mathpar}
    \end{minipage}
\end{center}

\paragraph{Event patterns}

The execution of \GRSync event patterns indicates if that event pattern is triggered or not, and
makes accessible the values of present events.
%
The initial states of a guarded event pattern \guarded{epat^+}{e_g} are denoted by
\initsync{epat^+}, and their transition with an evaluation context $\rho$ is denoted by
\begin{equation*}
    \epatOpsem{\rho}{s}{\guarded{epat^+}{e_g}}{s'}{b}{\rho'}.
\end{equation*}
The guard is optional and is only used in the rule for multiple event patterns. The transition
updates the internal state, produces an output bottom boolean $b_\bot$ indicating whether the event
is present and the guard holds ($\bot$ means it cannot yet be determined), and produces a new
evaluation context $\rho'$ that includes the values of identifiers corresponding to the present
events.
%
The initial state of guarded multiple event patterns (\guarded{epat^+}{e_g}) contains the initial
states of all event patterns ($epat$), as defined by \Cref{eq:semantics:grsync:opsem:init:epats}.
Its transition, ruled by \nameref{rule:semantics:grsync:opsem:step:epats:bool}, executes the
transitions of all individual event patterns ($epat$) and the transition of the guard $e_g$. The
event pattern holds only if all individual event patterns and the guard hold. The locally defined
identifiers are combined in a one evaluation context. If one event pattern returns $\bot$, then the
rule \nameref{rule:semantics:grsync:opsem:step:epats:bot1} applies and the tuple event pattern
returns $\bot$. If the guard returns $\bot$, then the rule
\nameref{rule:semantics:grsync:opsem:step:epats:bot2} applies and returns $\bot$ as well.

\begin{align}
    \stateTy{\guarded{epat^+}{e_g}}  & = \Pi_{epat} \stateTy{epat} & \text{\iff } & \stateTy{e_g} = \{()\} \label{eq:semantics:grsync:opsem:stateTy:epats} \\
    \initsync{\guarded{epat^+}{e_g}} & = (\initsync{epat})^+       & \text{\iff } & \initsync{e_g} = () \label{eq:semantics:grsync:opsem:init:epats}
\end{align}

\begin{center}
    \begin{minipage}{0.9\textwidth}
        \begin{mathpar}
            \inferOPSEMEPat
            {
                \left(\epatOpsem{\rho}{s}{epat}{s'}{b'}{\rho'}\right)^+ \\
                \exprOpsem{\rho \ctxconcat \rho'}{()}{e_g}{()}{b_g} \\
                b = b_g \kwf{\&\&} (\kwf{\&\&} b')^+
            }
            { \epatOpsem{\rho}{s^+}{\guarded{epat^+}{e_g}}{s^{\prime+}}{b}{\bigCtxconcat \rho^{\prime +}} }
            {Tple{\scriptsize bool}}
            \label{rule:semantics:grsync:opsem:step:epats:bool}

            \inferOPSEMEPat
            {
                \left(\epatOpsem{\rho}{s}{epat}{s'}{b'_\bot}{\rho'}\right)^+ \\
                \exists b_\bot \in \{b_\bot^{\prime+}\}, \: b_\bot = \bot
            }
            { \epatOpsem{\rho}{s^+}{\guarded{epat^+}{e_g}}{s^+}{\bot}{\bigCtxconcat \rho^{\prime +}} }
            {Tple{\scriptsize bot1}}
            \label{rule:semantics:grsync:opsem:step:epats:bot1}

            \inferOPSEMEPat
            {
                \left(\epatOpsem{\rho}{s}{epat}{s'}{b'}{\rho'}\right)^+ \\
                \exprOpsem{\rho \ctxconcat \rho'}{()}{e_g}{()}{\bot}
            }
            { \epatOpsem{\rho}{s^+}{\guarded{epat^+}{e_g}}{s^+}{\bot}{\bigCtxconcat \rho^{\prime +}} }
            {Tple{\scriptsize bot2}}
            \label{rule:semantics:grsync:opsem:step:epats:bot2}
        \end{mathpar}
    \end{minipage}
\end{center}

The event detection pattern (\detect{x}{y}) has an empty state, as
\Cref{eq:semantics:grsync:opsem:init:epat:event} states.
%
Its transition is determined by rules
\nameref{rule:semantics:grsync:opsem:step:epat:event:some},
\nameref{rule:semantics:grsync:opsem:step:epat:event:none}, and
\nameref{rule:semantics:grsync:opsem:step:epat:event:bot}, which check if the event $y$ is present
by looking at its associated value in the context $\rho$.
%
If $\rho[y] = \some{v}$, rule \nameref{rule:semantics:grsync:opsem:step:epat:event:some} applies:
the event $y$ is present, the pattern \detect{x}{y} holds, and its transition returns a context
containing the unwrapped value of the event ($[v/x]$).
%
Conversely, if the identifier $y$ is associated with a \none value, rule
\nameref{rule:semantics:grsync:opsem:step:epat:event:none} applies: the pattern \detect{x}{y} does not
hold, and its transition returns an empty context.
%
If the identifier $y$ is not determined (\ie associated with $\bot$), rule
\nameref{rule:semantics:grsync:opsem:step:epat:event:bot} applies: the pattern \detect{x}{y} cannot
yet be determined, it returns $\bot$ and an empty context.

\begin{align}
    \stateTy{\detect{x}{y}} & = \{()\} & \initsync{\detect{x}{y}} & = () \label{eq:semantics:grsync:opsem:init:epat:event}
\end{align}

\begin{center}
    \begin{minipage}{0.95\textwidth}
        \begin{mathpar}
            \inferOPSEMEPat
            { \none = \rho[y] }
            { \epatOpsem{\rho}{()}{\detect{x}{y}}{()}{\false}{\botrho{\{x\}}} }
            {Dtct{\scriptsize none}}
            \label{rule:semantics:grsync:opsem:step:epat:event:none}

            \inferOPSEMEPat
            { \some{v} = \rho[y] }
            { \epatOpsem{\rho}{()}{\detect{x}{y}}{()}{\true}{[v / x]} }
            {Dtct{\scriptsize some}}
            \label{rule:semantics:grsync:opsem:step:epat:event:some}

            \inferOPSEMEPat
            { \bot = \rho[y] }
            { \epatOpsem{\rho}{()}{\detect{x}{y}}{()}{\bot}{\botrho{\{x\}}} }
            {Dtct{\scriptsize bot}}
            \label{rule:semantics:grsync:opsem:step:epat:event:bot}
        \end{mathpar}
    \end{minipage}
\end{center}

The state of a rising edge $e$ contains the memory of the last evaluation of $e$, initialized to
\false by \Cref{eq:semantics:grsync:opsem:init:epat:edge}.
%
Its transition is described by the rule \nameref{rule:semantics:grsync:opsem:step:epat:edge:bool}:
it stores the previous boolean value of $e$ in its state and returns \true when it detects a
transition from \false to \true. If the transition of $e$ returns $\bot$, then the rising edge
cannot yet be determined and returns $\bot$.

\begin{align}
    \stateTy{e}  & = \Bool \times \stateTy{e}                                                        &
    \initsync{e} & = (\texttt{false}, \initsync{e}) \label{eq:semantics:grsync:opsem:init:epat:edge}
\end{align}

\begin{center}
    \begin{minipage}{0.9\textwidth}
        \begin{mathpar}
            \inferOPSEMEPat
            {
                \exprOpsem{\rho}{s}{e}{s'}{b'} \\
                b = \kwf{!}v_m \kwf{\&\&} b'
            }
            { \epatOpsem{\rho}{(v_m, s)}{e}{(b', s')}{b}{\emptyrho} }
            {Edge{\scriptsize bool}}
            \label{rule:semantics:grsync:opsem:step:epat:edge:bool}

            \inferOPSEMEPat
            { \exprOpsem{\rho}{s}{e}{s'}{\bot} }
            { \epatOpsem{\rho}{(v_m, s)}{e}{(v_m, s)}{\bot}{\emptyrho} }
            {Edge{\scriptsize bot}}
            \label{rule:semantics:grsync:opsem:step:epat:edge:bot}
        \end{mathpar}
    \end{minipage}
\end{center}

\paragraph{Summary}
This section presented the operational semantics of \GRSync using a small-step, state-machine-based
formalism. We introduced a set of transition rules that describe the execution of each language
construct, from high-level components to individual expressions and patterns. The state, which is
finite and statically determined, encapsulates memories and the states of nested components. A key
aspect of this semantics is the use of a fix-point iteration to evaluate simultaneous equations.
This approach resolves data dependencies without imposing a fixed execution order. The propagation
of a bottom value ($\bot$) ensures that computations proceed only when their inputs are determined.
